%% Section 7: Conclusion
%% Japanese draft (draft2/07_conclusion.md) embedded as comments.
%% English translation written below each paragraph.
%% -----------------------------------------------------------------------

\section{Conclusion}
\label{sec:conclusion}

% 本研究では IntSeqBERT を提案した。整数数列向けの双ストリーム Transformer エンコーダであり、各要素を
% 2 つの相補的な軸で表現する。1 つ目はスケールと成長挙動を捉える連続的な対数 Magnitude 埋め込みであり、
% 2 つ目は周期的算術構造を捉える Sin/Cos Modulo 埋め込みである。
% 2 つのストリームは FiLM で融合され、219,765 件の OEIS 数列に対して Magnitude 回帰・符号分類・
% 100 次元 Modulo 予測を組み合わせたマルチタスク目標で共同学習される。
We have presented \textbf{IntSeqBERT}, a dual-stream Transformer encoder that fuses a
continuous log-scale magnitude embedding with sin/cos modulo embeddings via FiLM, jointly
trained on 219,765 OEIS sequences. Experiments (Section~\ref{sec:experiments}) demonstrate
that the dual-stream design substantially outperforms both a standard tokenised Transformer
and a magnitude-only ablation across all scales and metrics.

% IntSeqBERT は Magnitude・符号・剰余を分離予測し CRT Solver で統合する設計であり、ニューラル
% ネットワークが学習しにくい厳密な数論的制約を記号処理が補う。純粋な end-to-end 学習の限界と
% ハイブリッド設計の必要性の両方を示す事例である。
Our results also carry a broader message. Even with number-theoretic structure built into
the representation, the absolute success rate of next-term prediction remains modest
(Top-1: 19.09\%), and a purely tokenised Transformer is far weaker still (2.59\%)---
confirming that neural networks struggle with direct prediction of OEIS terms.
IntSeqBERT itself is a neural-symbolic design: the network predicts magnitude, sign, and
residues separately, and the symbolic CRT solver supplies the exact number-theoretic
constraints that neural networks struggle to learn end-to-end. We therefore read our
results both as evidence of the limits of pure end-to-end learning on this task and as
motivation for hybrid symbolic--neural approaches to computer mathematics.

% 今後の課題として以下が挙げられる：
\textbf{Future work} includes combining magnitude uncertainty estimation with approximate CRT
to improve large-integer prediction; extending the modulo stream beyond $m = 101$;
evaluating on the remaining FACT benchmark tasks---sequence classification, sequence
similarity, next sequence-part prediction, and unmasking~\cite{zurich-fact}---to enable direct
comparison against that benchmark;
constructing family-aware dataset splits and introducing sequence-synthesis data
augmentation to improve generalisation in sparse magnitude buckets;
scaling to larger models pre-trained on the full OEIS corpus;
and exploring downstream applications such as conjecture generation,
where pretrained IntSeqBERT embeddings could serve as a rich feature backbone.
